\section{Frontend: Pruebas de Extremo a Extremo (E2E) y Accesibilidad}
\label{sec:s4_frontend}

El frontend entrega las vistas de visibilidad y privacidad del portafolio y se somete a un
control de calidad exhaustivo: \textit{cross-device}, flujos E2E y accesibilidad. La
Figura~\ref{fig:s4_visibilidad} muestra el panel de ajustes de visibilidad.\par

\begin{figure}[!ht]
    \centering
    \fbox{\begin{minipage}[c][4.0cm][c]{0.62\textwidth}
        \centering
        \sffamily\color{grisISO}\footnotesize
        [\,Captura de pantalla pendiente\,]\par
        \vspace{0.3cm}
        Ajustes de visibilidad, SEO y protección con contraseña del portafolio\par
        \vspace{0.2cm}
        \tiny(reemplazar por \texttt{assets/img/screenshots/visibilidad.png})
    \end{minipage}}
    \caption{Panel de visibilidad y privacidad del portafolio.}
    \label{fig:s4_visibilidad}
\end{figure}

\subsection{Inconsistencias Visuales y Cross-device Testing}
\label{ssec:s4_fe_crossdevice}
Se resolvieron inconsistencias visuales en distintas pantallas y resoluciones mediante
\textit{cross-device testing}, garantizando una adaptación fluida sin desbordamientos en móvil,
tableta y escritorio. La tabla~\ref{tab:s4_breakpoints} resume los \textit{breakpoints} verificados.

\begin{table}[!ht]
    \centering
    \renewcommand{\arraystretch}{1.3}
    \fontsize{9}{10.8}\selectfont\sffamily
    \begin{tabularx}{\textwidth}{|l|l|X|}
        \hline
        \rowcolor{colorPrimario}
        \textcolor{white}{\textbf{Dispositivo}} & \textcolor{white}{\textbf{Ancho}} & \textcolor{white}{\textbf{Verificación}} \\ \hline
        Móvil & 360--430\,px & Sidebar visible; sin scroll horizontal. \\ \hline
        \rowcolor{grisTecnico}
        Tableta & 768--1024\,px & Reflujo de columnas; tablas responsivas. \\ \hline
        Escritorio & $\geq$1280\,px & Sidebar oculto; layout completo. \\ \hline
    \end{tabularx}
    \caption{Breakpoints verificados en el cross-device testing (Sprint 4).}
    \label{tab:s4_breakpoints}
\end{table}

\subsection{Flujos Completos y Estándares de Accesibilidad}
\label{ssec:s4_fe_accesibilidad}
Se simularon automatizadamente flujos completos de navegación del perfil (Vitest + Testing
Library, entorno jsdom) y se verificaron los estándares de accesibilidad: contraste de color,
navegación por teclado y etiquetado semántico.

\begin{TablaCasoPrueba}
    1 & Navegación login-dashboard por teclado & Foco visible y orden lógico & \estadoPass \\ \hline
    2 & Contraste de texto sobre fondo primario & Cumple ratio AA (4.5:1) & \estadoPass \\ \hline
    3 & Render en viewport de 360px & Sin scroll horizontal ni cortes & \estadoPass \\ \hline
    4 & Portafolio protegido sin contraseña & Acceso bloqueado correctamente & \estadoPass \\ \hline
\end{TablaCasoPrueba}

\subsection{Herramientas y Criterios de Accesibilidad}
\label{ssec:s4_fe_a11y}
La verificación de accesibilidad se apoya en los criterios de la tabla~\ref{tab:s4_a11y},
alineados con las pautas WCAG nivel AA.

\begin{table}[!ht]
    \centering
    \renewcommand{\arraystretch}{1.3}
    \fontsize{9}{10.8}\selectfont\sffamily
    \begin{tabularx}{\textwidth}{|l|X|}
        \hline
        \rowcolor{colorPrimario}
        \textcolor{white}{\textbf{Criterio}} & \textcolor{white}{\textbf{Verificación}} \\ \hline
        Contraste & Ratio $\geq$ 4.5:1 en texto normal (AA). \\ \hline
        \rowcolor{grisTecnico}
        Navegación por teclado & Orden de foco lógico y foco visible. \\ \hline
        Etiquetado semántico & Roles ARIA y etiquetas en formularios. \\ \hline
        \rowcolor{grisTecnico}
        Movimiento reducido & Respeto de \comando{prefers-reduced-motion}. \\ \hline
        Responsividad & Sin scroll horizontal en 360--1280\,px. \\ \hline
    \end{tabularx}
    \caption{Criterios de accesibilidad verificados (Sprint 4).}
    \label{tab:s4_a11y}
\end{table}

Las pruebas automatizadas (Vitest + Testing Library, \comando{jsdom}) simulan flujos completos de
navegación; los hallazgos de contraste y foco se resolvieron antes del cierre del incremento,
conforme al DoD de la Sección~\ref{ssec:s0_dod}.

\clearpage
